\documentclass[aps,preprint]{revtex4}%
\usepackage{amssymb}
\usepackage{amsfonts}%
\usepackage{amsmath}%
\setcounter{MaxMatrixCols}{30}%
\usepackage{graphicx}
%TCIDATA{OutputFilter=latex2.dll}
%TCIDATA{Version=5.50.0.2960}
%TCIDATA{CSTFile=revtex4.cst}
%TCIDATA{Created=Monday, August 20, 2018 09:29:00}
%TCIDATA{LastRevised=Wednesday, August 22, 2018 08:55:12}
%TCIDATA{<META NAME="GraphicsSave" CONTENT="32">}
%TCIDATA{<META NAME="SaveForMode" CONTENT="1">}
%TCIDATA{BibliographyScheme=Manual}
%TCIDATA{<META NAME="DocumentShell" CONTENT="Articles\SW\REVTeX 4">}
%BeginMSIPreambleData
\providecommand{\U}[1]{\protect\rule{.1in}{.1in}}
%EndMSIPreambleData
\newtheorem{theorem}{Theorem}
\newtheorem{acknowledgement}[theorem]{Acknowledgement}
\newtheorem{algorithm}[theorem]{Algorithm}
\newtheorem{axiom}[theorem]{Axiom}
\newtheorem{claim}[theorem]{Claim}
\newtheorem{conclusion}[theorem]{Conclusion}
\newtheorem{condition}[theorem]{Condition}
\newtheorem{conjecture}[theorem]{Conjecture}
\newtheorem{corollary}[theorem]{Corollary}
\newtheorem{criterion}[theorem]{Criterion}
\newtheorem{definition}[theorem]{Definition}
\newtheorem{example}[theorem]{Example}
\newtheorem{exercise}[theorem]{Exercise}
\newtheorem{lemma}[theorem]{Lemma}
\newtheorem{notation}[theorem]{Notation}
\newtheorem{problem}[theorem]{Problem}
\newtheorem{proposition}[theorem]{Proposition}
\newtheorem{remark}[theorem]{Remark}
\newtheorem{solution}[theorem]{Solution}
\newtheorem{summary}[theorem]{Summary}
\newenvironment{proof}[1][Proof]{\noindent\textbf{#1.} }{\ \rule{0.5em}{0.5em}}
\begin{document}
\preprint{HEP/123-qed}
\title[Short title for running header]{REV\TeX4}
\author{First Author}
\affiliation{My Institution}
\author{Second Author}
\affiliation{My Institution}
\author{Third Author}
\affiliation{Other Institution}
\keywords{one two three}
\pacs{PACS number}

\begin{abstract}
Shell document for REV\TeX{} 4.

\end{abstract}
\volumeyear{year}
\volumenumber{number}
\issuenumber{number}
\eid{identifier}
\date[Date text]{date}
\received[Received text]{date}

\revised[Revised text]{date}

\accepted[Accepted text]{date}

\published[Published text]{date}

\startpage{101}
\endpage{102}
\maketitle
\tableofcontents


\section{AF and UHF $c^{\ast}$-algebras}

Bratelli defines \emph{Approximately Finite} dimensional (AF) $c^{\ast}%
$-algebras are the uniform closure of ascending sequences of finite
dimensional $c^{\ast}$-algebras (\emph{Inductive Limit}) i.e.
\begin{equation}
\mathcal{A}=u-\lim\limits_{r\rightarrow\theta}\mathcal{A}_{r},
\end{equation}
where
\begin{equation}
\mathcal{A}_{r}\subset\mathcal{A}_{r+1};1\leq r\leq\infty
\end{equation}
and $\mathcal{A}_{r}$ are $r$-dimensional $c^{\ast}$-algebras. 

Glimm defines $c^{\ast}$-algebras that are the uniform closure of strictly
ascending sequences of full $r\times r$ matrix algebras, all having the same
unit as \emph{Uniformly Hyperfinite} (UHF) algebras i.e
\begin{equation}
M\left(  \mathbb{C},\infty\right)  =u\text{-}\lim\limits_{r\rightarrow\theta
}M\left(  \mathbb{C},r\right)
\end{equation}


Dixmier removes the assumption that the matrix algebras have the same unit and
calls the resultant $c^{\ast}$-algebras \emph{Matroid }$c^{\ast}%
$\emph{-algebras}. The difference between AF algebras and the othe two is that
latter are simple algebras while the former can be non simple.  

\section{Ideals in $\mathcal{F}$}

$\mathfrak{I}_{L}$ is a left ideal if
\begin{equation}
\mathcal{F}\mathfrak{I}_{L}\subseteq\mathfrak{I}_{L}%
\end{equation}
and $\mathfrak{I}_{R}$ a right ideal if %

\begin{equation}
\mathfrak{I}_{R}\mathcal{F}\subseteq\mathfrak{I}_{R}%
\end{equation}
and a two sided ideal $\mathfrak{I}$ if
\begin{align}
\mathcal{F}\mathfrak{I}  & \subseteq\mathfrak{I}\\
\mathfrak{I}\mathcal{F}  & \subseteq\mathfrak{I.}%
\end{align}
An ideal $\mathfrak{I}$ is self adjoint Iff
\begin{equation}
A\in\mathfrak{I\Leftrightarrow}A^{\dagger}\in\mathfrak{I.}%
\end{equation}


\section{Inductive Limits}

Let $\left\{  \mathcal{A}_{n}\right\}  $ be a sequence of $c^{\ast}$-algebras
with injective $\ast$-homorphisms (injectiveness may not be necessary but it
is convenient)
\begin{equation}
\U{3c0} _{n}:\mathcal{A}_{n}\rightarrow\mathcal{A}_{n+1}.
\end{equation}
For $m>n$ we define
\begin{equation}
\U{3c0} _{nm}:\mathcal{A}_{n}\rightarrow\mathcal{A}_{m}.
\end{equation}
by%
\begin{equation}
\U{3c0} _{nm}=\U{3c0} _{m-1}\circ\U{3c0} _{m}\circ\cdots\circ\U{3c0}
_{n+1}\circ\U{3c0} _{n}%
\end{equation}
There is a unique (up to isomorphism) $c^{\ast}$-algebra $\mathcal{A}$ and
injective $\ast$-homomorphisms%
\begin{equation}
\U{3c3} _{n}:\mathcal{A}_{n}\rightarrow\mathcal{A},
\end{equation}
it is typical to think of $\mathcal{A}_{n}$ as subalgebras of $\mathcal{A}$
i.e. by identifying $\mathcal{A}_{n}$ with it's image in $\mathcal{A},$that
satisfies the following two conditions:

\begin{enumerate}
\item For $m>n$%
\begin{equation}
\U{3c3} _{n}=\U{3c0} _{nm}%
\end{equation}


\item The union
\begin{equation}%
%TCIMACRO{\tbigcup _{1\leq n\leq\infty}}%
%BeginExpansion
{\textstyle\bigcup_{1\leq n\leq\infty}}
%EndExpansion
\U{3c3} _{n}\left(  \mathcal{A}_{n}\right)
\end{equation}
is dense in $\mathcal{A}.$
\end{enumerate}

Let $\mathcal{A}$ $=(\mathcal{A}_{n},\U{3b1} _{n});n\in\mathbb{N}$, be an
inductive sequence of $c^{\ast}$-algebras then its inductive limit,
$\mathcal{A}$ is the enveloping $c^{\ast}$-algebra of the infinite union of
the $c^{\ast}$-algebras $\left\{  \mathcal{A}_{n}\right\}  $ and the $c^{\ast
}$-algebra $\mathcal{A}$
\begin{equation}
\mathcal{A}=\,\underset{n\rightarrow\infty}{i\text{-}\lim}\left\{
\mathcal{A}_{n},\U{3c0} _{n}\right\}
\end{equation}
is called the \emph{inductive limit} of the inductive sequence $\left\{
\mathcal{A}_{n},\U{3c0} _{n}\right\}  _{1\leq n\leq\infty}.$

. 

If $\mathcal{A}$ is an inductive limit of finite-dimensional $c^{\ast}%
$-algebras, we call $\mathcal{A}$ an AF-algebra, (Almost Finite). 

\bigskip






\end{document}